% critique.tex — f173v, vv. 920–958 : commentaire critique.
% Parties textuelles du commentaire de l'auteur (travaux antérieurs, Harmattan 2023), texte et notes
% conservés ; répartition, harmonisation et passages signalés : repartition.md.
\cc{943–56}{Le papyrus d'Oxyrhynque n° 3839\nc{\italique{The Oxyrhynchus Papyri}, vol. LVI (1989), n°3839, fr. 2. Analysé par W. E. H. Cockle avec le concours de Colin Austin.},}{dont il a été déjà question aux vers 742-766 contient également un fragment de ce passage. Il est si gâté qu'il ne nous apporte aucune information (si ce n'est une colométrie que Lætitia suivit \anglais{without strong conviction}\nc{L. P. E. Parker, \italique{The Songs of Aristophanes} (1997), p. 430.}). Le lecteur le trouvera néanmoins reproduit en annexe ; gage à lui d'en tirer quelque chose.}
\cc{947–1000}{}{Quant à sa structure, celle d'Enger, assez communément admise -- et notamment par Lætitia\nc{Dans son indispensable \italique{The Songs of Aristophanes}, pp. 428-436.} --, me semble juste  :\citComm{--- anapestes (v. 947-952) ;\\
--- proode (v. 953-958) ;\\
--- première paire strophique (v. 959-962 $=$ v. 963-965) ;\\
--- première mésode (v. 966-968) ;\\
--- seconde paire strophique (v. 969-976 $=$ v. 977-984) ;\\
--- seconde mésode (v. 985-989) ;\\
--- tierce paire strophique (v. 990-994 $=$ v. 995-1000).}}
